\section{研究基础与可行性分析}

\subsection{理论可行性分析}

研究点一所依赖的理论基础涵盖两个方面：一是访问控制与信息流约束，将"某一工具调用是否违反任务契约的信息流规则"表述为在统一执行轨迹图上，对来源节点到受保护对象的路径查询与契约条件判定，属于访问控制和信息流分析领域中具有长期研究积累的确定性判定范式；二是程序分析与图路径检查，将运行时事件、工具调用参数和返回值之间的依赖关系组织为区分证据等级的有向异构图，并在图上执行SAST风格的增量规则检查，这类方法在系统安全领域的溯源图分析和静态分析中均有成熟的算法基础。研究点二所依赖的理论基础是图表示学习与边级分类，将Runtime前缀图中已观测边的安全角色定位（SOURCE/SINK/BENIGN）建模为图神经网络与预训练语言模型联合驱动的边级分类与匹配排序问题，该范式在知识图谱补全、链路预测和图异常检测等领域已有广泛的方法支撑。需要说明的是，将上述成熟范式迁移到LLM Agent工具调用链安全场景均需要面对"运行时事件的异构语义""图结构需适配Agent执行语义""边标签需基于benchmark强监督可重复生成"等具体适配问题，前人在相邻范式上的探索为本研究提供了方法论参照，但并不能替代本研究在该具体场景下的适配与实证工作。

\subsection{实验可行性分析}

实验条件方面，本研究依托的AgentDyn等开源Agent安全评测脚手架已提供完善的工具执行环境、参数规范化流程与任务效用/攻击成功评测器，具备支撑Agent开发与安全评测实验的软硬件基础，可支撑课题的主实验、消融实验与统计检验方案的实施。数据可行性方面，AgentDyn基准在AgentDojo已有的银行、工作区、旅行与Slack四个静态场景基础上，扩展了购物、代码协作与日常生活三个动态开放场景，其任务设计强调动态变化的用户状态与工具返回内容，攻击载荷通过邮件、商品评论、代码托管平台评论等多种真实数据通道注入，与本研究面向的"跨步骤间接提示注入"场景高度契合，能够为研究点一的证据图构建与规则检查以及研究点二的攻击边定位模型训练与评测分别提供充足的正常任务与攻击任务样本。实验设计模式方面，本研究采用的"以评测套件加用户任务为独立分析单元、配对bootstrap置信区间估计效应量、多重比较采用Holm校正符号检验"的统计检验方案，是安全评测领域已被广泛验证的成熟范式，评估指标（安全任务完成率、干预范围、排序准确性、检测与恢复相关指标、成本指标）均可由评测框架的执行日志直接、客观计算，不依赖人工主观打分，具备指标客观可信、实验流程可复现的特点。

\subsection{已有研究基础}

文献调研与课题查新方面，已围绕LLM Agent安全、间接提示注入攻防、执行溯源与图表示学习等方向，完成2021--2026年相关文献的系统调研，覆盖提示工程类、训练检测类与架构设计类三条主要防御技术路线，并结合最新的证据地图对邻近工作的方法边界与局限进行了逐篇核验，为课题的新颖性与原创性提供了初步验证依据。实验工作方面，已在AgentDojo、AgentDyn两个数据集上对多种防御基线方法完成初步适配，搭建了实验脚手架并进行了前期验证，形成了一批用于检验实验流程可行性的初步对比分析数据。
